\documentclass{article}
\usepackage[UTF8]{ctex}
\usepackage{amsmath,amssymb,amsfonts}
\usepackage{enumitem}
\usepackage{geometry}
\geometry{a4paper, margin=1in}
% === 错题本排版优化（不参与题目标识） ===
\setlength{\parindent}{0pt}
\setlength{\parskip}{0.35em plus 0.10em minus 0.08em}
\setlength{\emergencystretch}{3em}
\setlist{itemsep=0.25em, topsep=0.40em, parsep=0pt}
\allowdisplaybreaks[2]
% === 错题本排版优化结束 ===
\begin{document}
\section*{吉林大学2024年数学分析真题}
\section*{75．吉林大学2024年数学分析试卷}

\subsection{题目1}
1． \(\lim _{n \rightarrow \infty} \frac{1+\frac{1}{2}+\cdots+\frac{1}{n^{2}-1}+\frac{1}{n^{2}}}{n \ln n}\) ．

\subsection{题目2}
2． \(\lim _{n \rightarrow \infty} n^{2}\left(n \sin \frac{1}{n}-1\right)\) ．

\subsection{题目3}
3． \(\lim _{x \rightarrow 0} \frac{\sqrt[3]{1+\tan ^{2} x}-\sqrt{\cos x}}{x \sin x}\) ．

\subsection{题目4}
4． \(\lim _{n \rightarrow \infty} n \int_{0}^{1} x^{n} \cos \left(1+x^{2}\right) \mathrm{d} x\) ．

\subsection{题目5}
5． \(\lim _{n \rightarrow \infty} \frac{1}{n}\left(\frac{1}{\sqrt{n^{2}+1}}+\frac{1}{\sqrt{n^{2}+2}}+\cdots+\frac{1}{\sqrt{2 n^{2}-1}}+\frac{1}{\sqrt{2 n^{2}}}\right)\) ．

\subsection{题目6}
6．设 \(u=u(x, y), v=v(x, y)\) 是由方程组

\[
\left\{\begin{array}{l}
x u-y v=0 \\
y u+x v=2
\end{array}\right.
\]

在点 \(\left(x_{0}, y_{0}, u_{0}, v_{0}\right)=(1,1,1,1)\) 的某个邻域内定义的隐函数，求 \(\frac{\partial u}{\partial x}(1,1), \frac{\partial u}{\partial y}(1,1), \frac{\partial v}{\partial x}(1,1), \frac{\partial v}{\partial y}(1\)

\subsection{题目7}
7．求曲面 \(2 z=x^{2}+y^{2}\) 与 \(z=\sqrt{x^{2}+y^{2}}\) 围成区域的体积．

\subsection{题目8}
8．求第二型曲线积分 \(\int_{L} \frac{x \mathrm{~d} y-y \mathrm{~d} x}{2 x^{2}+y^{2}}\) ，其中 \(L\) 是圆周 \(x^{2}+y^{2}=1\) ，方向为逆时针．

二．求 \(f(x, y)=2 x^{2}+4 x y-2 y^{2}\) 在闭区域 \(D=\left\{(x, y) \mid x^{2}+y^{2} \leq 5\right\}\) 上的最大值和最小值．

三．求级数 \(\sum_{n=1}^{\infty} \frac{(-1)^{n-1}}{n 2^{n}}\) 的和．

四．设 \(a_{n}=\left(1-\frac{2 \ln (\ln n)}{n}\right)^{n}, n=2,3, \cdots\), 判断级数 \(\sum_{n=2}^{\infty} a_{n}\) 的敛散性．

五．证明： \(\lim _{h \rightarrow 0^{+}} \int_{0}^{1} \frac{h}{h^{2}+x^{2}} \ln \left(x^{2}+2\right) \mathrm{d} x=\frac{\pi}{2} \ln 2\) ．

六．设函数 \(f(x)\) 在 \([0,2]\) 上二阶导数连续，且 \(f(0)=f(2)=0\) ，证明：

\[
\int_{0}^{2} f(x) \mathrm{d} x \leq \frac{2}{3} \max _{x \in[0,2]}\left\{\left|f^{\prime \prime}(x)\right|\right\}
\]

七．设 \(F\) 为定义在 \(\mathbb{R}^{n}\) 上的函数．\\
（1）函数 \(F\) 不一致连续的充要必要条件是：存在点列 \(\left\{\mathbf{x}^{(n)}\right\}\) 和 \(\left\{\mathbf{y}^{(n)}\right\} \subset \mathbb{R}^{n}\) 和常数 \(a>0\) ，其中有\(\lim _{n \rightarrow \infty}\left\|\mathbf{x}^{(n)}-\mathbf{y}^{(n)}\right\|=0\) ，但对 \(n=1,2, \cdots\) ，都有 \(\left|F\left(\mathbf{x}^{(n)}\right)-F\left(\mathbf{y}^{(n)}\right)\right| \geq a\) 。\\
（2）若 \(F\left(x_{1}, x_{2}, \cdots, x_{n}\right)=f\left(x_{1}^{2}+x_{2}^{2}+\cdots+x_{n}^{2}\right), f\) 在 \([0,+\infty)\) 上连续可导，同时 \(\lim _{t \rightarrow+\infty} f^{\prime}(t)=b \neq 0\)证明：\(F\) 在 \(\mathbb{R}^{n}\) 上不一致连续．
\end{document}
